% !PRINTED 585
\weberpage{585}

We wish here also to emphasise the theorem expressed in
formula (34):

\gesperrt{Every prime of the form} $4f+1$ \gesperrt{can be decomposed into
the sum of two squares.}

\begin{center}
\S.~173.\\[6pt]
\gesperrt{The complex numbers of Gauss.}
\end{center}

The theorem last proved forms the foundation for the
theory of the number field which arises from the field of rational
numbers by adjunction of the imaginary
unit $i=\sqrt{-1}$, and which by our stipulation we have to denote
by $R(i)$\textsuperscript{1}).

We have here to distinguish the primes of the form $4f+1$
from those of the form $4f+3$, and for the sake of brevity we
wish to combine the former with the prime 2 and denote them
by $p$, the latter by $q$; thus we set
\[
p=4f+1,2;\quad q=4f+3
\]
The following theorems then hold:

1. \gesperrt{For every $p$ there are two integers $a$ and $b$
such that}
\[
p=a^2+b^2
\]
\gesperrt{holds.}

For the primes of the form $4f+1$ this has been proved in the
concluding theorem of the last section, and for $p=2$ it is
immediately evident from $2=1^2+1^2$.

Against this stands a second theorem:

2. \gesperrt{A prime $q$ is never representable in the form
$a^2+b^2$,} or the still more general one:

3. \gesperrt{The sum of two squares of integers $a^2+b^2$
is divisible by a prime $q$ only when $a$ and $b$
are divisible by $q$.}

Theorem 2. follows simply from the fact that, when $a^2+b^2$
is odd, one of the two numbers $a$, $b$ is even, and hence its
square divisible by 4, while the other is odd, and hence its square
%CONT>>

\textsuperscript{1}) \gesperrt{Gauss}, Theoria residuorum biquadraticorum, commentatio
secunda. Werke, Bd.~II.
